% core_thesis.tex

\cleardoublepage
\chapter*{Core Thesis}
\phantomsection
\addcontentsline{toc}{chapter}{Core Thesis}

\begin{center}
\large
\textit{
Alignment is not a curve.\\
It is a curvature-driven, realization-aware generative system.
}
\end{center}

\vspace{1em}

\begin{remark}
Classical railway geometry represents track layouts as curves in
Euclidean space.
While sufficient for visualization, this perspective is inadequate for
design, construction, and dynamic analysis.
\end{remark}

\begin{remark}
Within the AIM framework, the fundamental object is not the curve
\(\gamma(s)\), but the curvature law \(\kappa(s)\), complemented by the
cant function \(\phi(s)\).
All geometric quantities are derived by integration.
\end{remark}

\begin{remark}
This shifts the modelling paradigm from descriptive geometry to
generative structure:
alignments are defined by the laws that produce curves, not by the
curves themselves.
\end{remark}

\begin{remark}
The overall AIM interpretation is summarized schematically in
\cref{fig:aim-overview}.
\end{remark}

\section*{Key Principles}

\begin{description}

\item[Curvature as Primary Variable]
The alignment is defined by curvature \(\kappa(s)\).
Position \(\gamma(s)\) is a derived quantity.
This inversion simplifies modelling and exposes the intrinsic structure
of railway geometry.

\item[Realization as Operator]
The physically constructed track differs from the theoretical alignment.
This transformation is captured by a realization operator
\(\mathcal{R}\), which includes discretization, interpolation, and
mechanical response.
Design must therefore be formulated in realized space.

\item[Hybrid Representation]
Real-world data and engineering design require both structure and
flexibility.
A hybrid model combines parametric representations with local
corrections, enabling robust reconstruction from imperfect data.

\item[Analytical Structure]
All relevant quantities arise from arc-length based relations.
As a consequence, derivatives with respect to parameters are available
in closed form.
This eliminates the need for numerical differentiation.

\item[Efficient Optimization]
Alignment optimization reduces to a structured nonlinear least-squares
problem with sparse and banded Jacobians.
This enables efficient solution via Gauss--Newton or SQP methods.

\end{description}

\section*{Implication}

\begin{remark}
The AIM framework unifies geometry, dynamics, realization, and data
processing within a single coherent model.

It replaces curve-based thinking by a system-based perspective in which
alignment, realization, and optimization are intrinsically linked.
\end{remark}

\begin{remark}
This provides a foundation for alignment-based information modelling,
capable of bridging theoretical design and practical railway
engineering.
\end{remark}
